\documentclass[12pt, a4]{article}

\usepackage{amsmath}
\usepackage{german}
\usepackage{graphicx}

%\graphicspath{{/home/ritz/Teaching/fcn/fig/}}

%----------------- layout ----------------------------------------

\oddsidemargin=0in
\textwidth=6in
\topmargin=-0.7in
\textheight=9.5in
\baselineskip=16pt
\font\mittel=cmr17
\font\roman=cmr12
\font\klein=cmr9
\font\mgross=cmr10
\newfont{\CAP}{cmcsc10 scaled 900}
\pagestyle{empty}
\parindent=0mm

%----------------- macros ----------------------------------------
\newcommand{\conc}[1]{[{\rm C}]_{\text{#1}}}
\newcommand{\units}[1]{\text{#1}}
\newcommand{\kalium}{$\text{K}^+$}
\newcommand{\natrium}{$\text{Na}^+$}
\newcommand{\chlor}{$\text{Cl}^-$}
\newcommand{\wasser}{$\text{H}_2\text{O}$}
%-----------------------------------------------------------------


\begin{document}

\parbox{2cm}{
\includegraphics[width=1.8cm]{bccnlogo}
}
\parbox{11cm}{
\begin{center}
\large HUMBOLDT-UNIVERSIT"AT \hskip 0.1 cm ZU \hskip 0.1 cm BERLIN
\vskip 0.1 true cm
\mgross BERNSTEIN CENTER FOR COMPUTATIONAL NEUROSCIENCE
\end{center}
}
\parbox{2cm}
{
\hfill
\includegraphics[width=1.8cm]{hublogo}
}

\vskip 0.6 true cm

\leftline{\CAP Humboldt-Universit"at zu Berlin 
\hfill Phone: 030/2093-9110}
\leftline{\CAP Philippstr. 13 House 6
              \hfill Fax: 030/2093-6771}
\leftline{\CAP \hfill webpage: http://www.bccn-berlin.de/}




\vskip 0.8 true cm
\centerline{\bf Models of Neural Systems I, WS 2010/11}
\centerline{\bf Computer Practical 2}
\centerline{Solutions to hand in on: November, 8th, 2010}


\vskip 0.8 true cm

\begin{enumerate}

    \item \textbf{Python modules}

    Python functions are usually organized into modules, which eases a lot the managing of projects. Here you will learn how to write and use simple modules.
    \begin{enumerate}
        \item Write your Fibonacci function into a file \texttt{mymodule.py} (you can also choose another name).
        \item Import \texttt{mymodule} module into the \texttt{ipython} shell and run your \texttt{fib} function from the previous sheet:
        \begin{verbatim}
            import mymodule
            mymodule.fib(10)
        \end{verbatim} 
        \item Try alternative forms to \texttt{import} functions:
        \begin{verbatim}
            import mymodule as my
            from mymodule import fib
            from mymodule import *
            reload(my)
        \end{verbatim}
        Investigate the differences between these cases.
        \item When you work in the interactive \texttt{ipython} shell, also the magic function \texttt{\%run} may prove useful. Find out from the documentation what the function does and use it to import your Fibonacci function.
        \item (Optional) It is also possible to write a \texttt{main} function, which will be called when the module is run as a script (for example like that: \texttt{python mymodule.py}). The simplest function looks like that:
        \begin{verbatim}
            ...
            def main():
                fib_series = mymodule.fib(10)
                print "Fibonacci series:", fib_series
            if __name__ == '__main__':
                main()
        \end{verbatim}
    \end{enumerate}

    \item \textbf{Numerical Python} (NumPy) \\
    The standard version of Python does not provide any functions for efficient numerical computation. However, the \texttt{scipy} module extends its functionality into this area. \texttt{Scipy} includes optimized algorithms for the creation and manipulation of numerical arrays, basic linear algebra operations, and others.
    
    \begin{enumerate}
        \item Import \texttt{numpy} (for example, \texttt{import numpy as np})
        \item Create the following arrays: integers from 0 to 99 (\texttt{np.arange}), the sequence $\lbrace5.1, 5.2, 5.3, ...,15.0\rbrace$ (\texttt{np.arange}), the 5x10 matrix of zeros or ones (\texttt{np.zeros}/\texttt{np.ones}), and the 5x5 identity matrix (\texttt{np.eye})
        \item Define two 2D arrays:
        $$
        A=\left(
        \begin{array}{ccc}
        1 & 2 & 3 \\
        2 & 1 & 2 \\
        3 & 2 & 1 \\
        \end{array}
        \right) \;
        B=\left(
        \begin{array}{ccc}
        0 & 2 & 1 \\
        3 & 0 & 2 \\
        1 & 3 & 0 \\
        \end{array}
        \right)
        $$
        and
        \begin{enumerate}
            \item multiply the array $A$ by a scalar
            \item multiply arrays $A$ and $B$ element-wise
            \item check if matrices $A$ and $B$ commute under the matrix dot product (use \texttt{np.dot} and $AB\stackrel{!}{=}BA$)
            \item Multiply element-wise column 1 of matrix A with row 2 of matrix B using the slice notation.
        \end{enumerate}
    \end{enumerate}

    \item \textbf{Basic plotting} (Pylab)
    \label{ex:pylab}
    
    Pylab (aka matplotlib) is a Python plotting library. Its  capabilities and interface are quite similar to the ones offered by Matlab. It can be used  to display the results of your calculations, create publication quality figures and can be embedded in graphical user interfaces.
    
    A simple plot in pylab can be drawn with the following example (taken from matplotlib tutorial):
    \begin{verbatim}
        import pylab as plt
        plt.figure()
        plt.plot([1,2,3,4])
        plt.show()
    \end{verbatim} 

    \begin{enumerate}
        \item Modify the above code, so that only points are drawn (without connecting lines).
        \item Label the x-axis, the y-axis and add a title to the figure (\texttt{plt.xlabel}, \texttt{plt.ylabel}, \texttt{plt.title}).
        \item Split the figure vertically into two panels (\texttt{plt.subplot}). In the upper panel, plot the data with lines. In the lower panel, plot it with bars (\texttt{plt.bar}).
        \item Save the figure in your home directory (you can use an icon from the figure's toolbar or the \texttt{plt.savefig} function)
    \end{enumerate}

    \item \textbf{Activation functions}
    \label{ex:activation}

        Plot the following functions in one plot. Add the axes' labels and a legend to the plot. Give an interpretation for the parameters of the functions.
    
        \begin{enumerate}
            \item Sigmoid function:
                $$f(x)=\frac{1}{1+\exp(-ax)}$$
                where $a$ is a free parameter.
            \item Hyperbolic tangent function:
                $$g(x)=\tanh(x).$$
            \item Threshold function:
                \begin{equation*}
                    h(x)=\begin{cases}1 & \text{if $x\geq0$} \\
                                    0 & \text{if $x<0$.}
                        \end{cases}
                \end{equation*}
            \item Piecewise linear function:
                \begin{equation*}
                    i(x)=\begin{cases} \frac{1}{2}& \text{if $x\geq\frac{1}{2}$} \\
                        x & \text{if $\frac{1}{2}>x>-\frac{1}{2}$} \\
                        -\frac{1}{2} & \text{if $x\leq-\frac{1}{2}$.}
                        \end{cases}
                \end{equation*}
    
        \end{enumerate}
    
        Export the figure to a file.

    \item \textbf {Report}

    Generate a single PDF file including your name(s), student identification number(s), figures from Exercises \ref{ex:pylab} and \ref{ex:activation}, captions, a short comment to your solutions and a full source code. To do that you can use  \texttt{pyreport} (see the course homepage), LaTeX or any other word processing program with PDF export. State precisely with each answer, which exercise it is related to (you can also copy the text of the exercise at the top of your solution). Please attach the source code in raw text format as well (as \texttt{*.py} file).
    
    \textbf{Send the report to your tutor not later than Monday, Nov 8nd, 9am!} 
\end{enumerate}

\vfill
\centerline{\CAP Contact}
\CAP

\begin{tabular}{lll}
Urs Bergmann & Phone: 2093-8924 & Email: ubergmann(at)fias.uni-frankfurt.de \\
Matthias Guggenmos & & Email: matthias.guggenmos(at)bccn-berlin.de \\
Richard Kempter \hfill & Phone: 2093-8925 \hfill & Email: r.kempter(at)biologie.hu-berlin.de \\
Robert Schmidt & Phone: 2093-8926 & Email: r.schmidt(at)biologie.hu-berlin.de \\
Bartosz Telenczuk & Phone: 2093-8838 & Email: b.telenczuk(at)biologie.hu-berlin.de \\
\end{tabular}

\end{document}
 
